\section*{Bab \ref{ch:b3:endocrinology} --- Endokrinologi dan Homeostasis}

\begin{solution}{exo:b3:endocrinology:1}
\omterm{def:b3:endocrinology:glucose}{Insulin}: peptida, reseptor permukaan (sebuah tirosin kinase), bekerja dalam
hitungan menit, paruh hidup sekitar \qty{5}{min}. Kortisol: steroid, \omterm{def:b3:endocrinology:hormone}{reseptor
inti}, berjam-jam, paruh hidup \qty{80}{min} (kebanyakan terikat pada sebuah pengusung).
Adrenalin: amina, reseptor tergandeng-protein~G di permukaan, hitungan detik,
paruh hidup sekitar \qty{2}{min}. Tiroksin: asam amino beriodin, \omterm{def:b3:endocrinology:hormone}{reseptor
inti}, berhari-hari, paruh hidup \qty{7}{d} (terikat pada pengusung).
Vasopresin: peptida, reseptor tergandeng-protein~G di permukaan (cAMP di dalam
\omterm{def:b3:renal-osmoregulation:nephron}{saluran pengumpulnya}), hitungan menit, paruh hidup sekitar \qty{15}{min}.
\end{solution}

\begin{solution}{exo:b3:endocrinology:2}
\omterm{def:b3:endocrinology:axes}{Hipotalamus} (TRH) $\to$ \omterm{def:b3:endocrinology:axes}{hipofisis} (TSH) $\to$ tiroid (T$_{4}$), dengan
T$_{4}$ menghambat kedua tingkat atasnya. Tiroid yang musnah: T$_{4}$ rendah,
TSH tinggi. \omterm{def:b3:cancer-biology:hallmarks}{Tumor} yang menyekresikan TSH: TSH tinggi dan T$_{4}$ tinggi (umpan
baliknya gagal menekan \omterm{def:b3:cancer-biology:hallmarks}{tumornya}). Kelebihan tablet tiroksin:
T$_{4}$ tinggi, TSH tertekan. Kekurangan iodin: T$_{4}$ rendah atau
rendah-normal, TSH tinggi, dan kelenjarnya tumbuh di bawahnya menjadi sebuah gondok.
\end{solution}

\begin{solution}{exo:b3:endocrinology:3}
Otaknya mati karena glukosa yang rendah dalam hitungan menit lalu dirugikan glukosa yang tinggi
hanya bertahun-tahun; dan evolusi menjumpai kelaparan jauh lebih sering daripada pesta.
Jadi pertahanan terhadap sebuah penurunan berlebihan (\omterm{def:b3:endocrinology:glucose}{glukagon}, adrenalin,
kortisol, \omterm{def:b3:endocrinology:hormone}{hormon} pertumbuhan) dan pertahanan terhadap sebuah kenaikan tunggal.
Ramalannya: \omterm{def:b3:endocrinology:glucose}{insulin} yang terlalu banyak mematikan secara akut (koma hipoglikemik),
yang terlalu sedikit adalah sebuah penyakit menahun --- dan itulah yang dilihat kliniknya.
\end{solution}

\begin{solution}{exo:b3:endocrinology:4}
Sebuah \omterm{def:b3:endocrinology:clock}{zeitgeber} adalah sebuah isyarat lingkungan yang menambat sebuah jam: cahaya,
waktu makan, suhu, olahraga, jadwal sosial. Jam induknya
memakai cahaya, lewat \omterm{def:b3:sensory-systems:retina}{sel ganglion} melanopsin \omterm{def:b3:sensory-systems:retina}{retinanya}. Sesudah
delapan zona waktu, jamnya bergeser hanya sekitar sejam tiap hari, sehingga selama
sepekan tidur, kortisol, suhu, dan seleranya berjalan pada waktu lamanya
sementara jam tepinya, yang disetel ulang santapannya, hanyut pada lajunya sendiri:
itulah lelah terbang, yang lebih buruk ke arah timur, sebab memajukan jamnya lebih sukar daripada
menundanya.
\end{solution}

\begin{solution}{exo:b3:endocrinology:5}
Tanggapan setengah maksimalnya pada $250$ reseptor yang terhuni: $H = K_{d}\times
250/(\num{20000} - 250) = \qty{1}{nM}\times 0.0127 = \qty{12.7}{pM}$,
yakni delapan puluh kali di bawah $K_{d}$. Dengan \num{2000} reseptor: $250/1750
\times \qty{1}{nM} = \qty{143}{pM}$, yakni sebelas kali kurang peka.
\omterm{def:b3:endocrinology:glucose}{Insulin} tinggi yang lama menurunkan pengaturan reseptornya, sehingga sebuah pengaruh tertentu
memerlukan lebih banyak \omterm{def:b3:endocrinology:glucose}{insulin}: yakni satu komponen kekebalan \omterm{def:b3:endocrinology:glucose}{insulin}, yang memberi makan
dirinya sendiri.
\end{solution}

\begin{solution}{exo:b3:endocrinology:6}
$T_{4} = 13$: $1.5\,\mathrm{e}^{0.92} = \qty{3.8}{mU/L}$; $11$:
$1.5\,\mathrm{e}^{1.84} = \qty{9.4}{mU/L}$; $9$: $1.5\,\mathrm{e}^{2.76}
= \qty{24}{mU/L}$; $25$: $1.5\,\mathrm{e}^{-4.6} = \qty{0.015}{mU/L}$. Pada
\qty{11}{pmol/L} T$_{4}$-nya masih di dalam \qtyrange{10}{20}{} sementara
TSH-nya lebih daripada dua kali batas atasnya (hipotiroidisme subklinis);
pada $13$ TSH-nya pas di dalam jangkauannya; pada $9$ keduanya tak normal; pada $25$
keduanya tak normal ke arah sebaliknya.
\end{solution}

\begin{solution}{exo:b3:endocrinology:7}
$L = s\beta/(a\gamma) = 0.2\times 0.05/(0.02\times 0.1) = 5$. Kelebihan
mantapnya $g^{*} = (u/a)/(1 + L) = (2/0.02)/6 = \qty{16.7}{mg/dL}$; tanpa
umpan balik $u/a = \qty{100}{mg/dL}$. Kelebihan \omterm{def:b3:endocrinology:glucose}{insulinnya} $i^{*} = \beta g^{*}/
\gamma = 0.05\times 16.7/0.1 = \qty{8.3}{mU/L}$.
\end{solution}

\begin{solution}{exo:b3:endocrinology:8}
Diskriminannya $(a - \gamma)^{2} - 4s\beta = 0.0064 - 0.04 < 0$:
berayun. $\omega = \sqrt{a\gamma + s\beta - (a+\gamma)^{2}/4} =
\sqrt{0.002 + 0.01 - 0.0036} = \qty{0.092}{min^{-1}}$, kalanya $2\pi/
\omega = \qty{68}{min}$; amplitudonya turun sebesar $\mathrm{e}$ dalam $2/(a +
\gamma) = \qty{16.7}{min}$. Dengan membagi dua $s$: $s\beta = 0.005$, diskriminannya
$0.0064 - 0.02 < 0$, masih berayun, $\omega = \sqrt{0.007 - 0.0036}
= \qty{0.058}{min^{-1}}$, kalanya \qty{108}{min}; waktu peluruhannya
tidak berubah; $L = 2.5$. Sebuah gelung yang lebih lemah kembali lebih lambat lalu menahan sebuah
galat mantap yang lebih besar.
\end{solution}

\begin{solution}{exo:b3:endocrinology:9}
Selama setahun prednisolonnya menekan CRH dan ACTH-nya; korteks
adrenalnya, yang tak terangsang, mengecil. Ketika berhenti, tidak ada steroid dari
luar, \omterm{def:b3:endocrinology:axes}{hipotalamus} dan \omterm{def:b3:endocrinology:axes}{hipofisisnya} memakan waktu berminggu-minggu untuk melanjutkan, dan
adrenalnya yang menyusut tidak dapat menanggapi ACTH bila ACTH datang. Sebuah infeksi
menuntut sepuluh kali kortisol normalnya; tanpa satu pun, pembuluhnya melebar,
tekanan dan glukosanya turun: itulah krisis adrenal. Yang terukur: ACTH rendah (lesinya
di pusat), kortisol rendah, dan kenaikan kortisol yang buruk terhadap ACTH yang
disuntikkan (kelenjarnya telah menyusut) --- tidak seperti penyakit Addison, yang ACTH-nya
tinggi. Karena itulah penurunan dosisnya perlahan.
\end{solution}

\begin{solution}{exo:b3:endocrinology:10}
Satu peubah: $\dot x = f(x)$ dengan $f' < 0$ punya sebuah titik tetap tunggal
$x^{*}$; bagi $x > x^{*}$, $\dot x < 0$ dan $x$ menurun secara monoton
ke arah $x^{*}$, tanpa pernah melintasinya (ketunggalan penyelesaiannya), dan
setangkup di bawahnya: tanpa ayunan. Dua peubah: matriksnya dapat punya
nilai eigen kompleks, seperti pada teoremanya, ketika pelaksananya bekerja lewat
sebuah peubah kedua dengan skala waktunya sendiri. Sebuah penekanan seketika
akan menetap pada sebuah aras yang mantap; penundaan antara transkripsi dan
penekanan intinya (terjemahan, pendimeran, fosforilasi, pengimporan)
membuat proteinnya menembus ke atas sebelum gennya dimatikan lalu menembus ke bawah
sebelum gennya dinyalakan lagi, dan kalanya kira-kira dua kali jumlah
penundaannya dan waktu merombak proteinnya. Perombakan PER yang lebih cepat
memendekkan kalanya: sebuah mutasi PER2 manusia yang mengacaukan kemantapannya memberi sebuah
jam sekitar 20 jam dan sebuah keluarga yang tertidur pukul 7 malam.
\end{solution}

\begin{solution}{exo:b3:endocrinology:11}
Pangkatnya $k(\text{Ca} - 1.2)$: pada $1.10$, $\mathrm{e}^{-2}$, PTH $= 1/
(1 + 0.135) = 0.88$; $1.15$: $0.73$; $1.20$: $0.50$; $1.25$: $1/(1 +
2.72) = 0.27$; $1.30$: $1/(1 + 7.39) = 0.12$. Kemiringan pada \omterm{thm:b3:endocrinology:loop}{titik setelnya}
$-k/4 = -5$ tiap mmol/L: yakni \qty{5}{\%} maksimumnya tiap
\qty{0.01}{mmol/L}. Kalsiumnya harus ditahan dalam beberapa persen dan
pelaksananya bekerja pada sebuah penyangga raksasa (tulangnya) tanpa tembusan atas, sehingga sebuah
tanggapan yang curam aman dan diperlukan. Sebuah gelung glukosa yang securam itu, yang bekerja lewat
\omterm{def:b3:endocrinology:glucose}{insulin} dengan waktu pembersihannya sendiri, akan punya \omterm{thm:b3:endocrinology:loop}{bati gelung} yang besar lalu
berayun dalam-dalam ke dalam hipoglikemia sesudah tiap santapan.
\end{solution}

\begin{solution}{exo:b3:endocrinology:12}
Glikasinya tak terpulihkan dan lambat, sehingga sebuah sel berumur $\tau$ punya
sebuah pecahan $kG\tau$ yang terglikasi; bila dirata-rata atas sel yang berumur seragam
\numrange{0}{120} hari, HbA1c $= kG\times\qty{60}{d}$, yakni sebanding
dengan glukosa rata-ratanya. Peneraannya: $5.5\times 60k = 0.05$ memberi $k =
\qty{1.5e-4}{}$ tiap (mmol/L)(hari); pada \qty{10}{mmol/L}: \qty{9.1}{\%}.
Dengan sel yang hidup 60 hari, umur rata-ratanya 30 hari dan HbA1c-nya
terbelah dua bagi glukosa yang sama: seorang pasien pada \qty{10}{mmol/L} akan terbaca
\qty{4.5}{\%}, yakni normal.
\end{solution}

\begin{solution}{pb:b3:endocrinology:1}
\textbf{1.} \qty{90}{mg/dL} $= \qty{0.9}{g/L} = 0.9/180 = \qty{5.0}{mmol/L}$;
$0.9\times 15 = \qty{13.5}{g}$ di dalam volume sebarannya.
\textbf{2.} $75/15 = \qty{5}{g/L}$: yakni sebuah kenaikan \qty{500}{mg/dL},
\qty{28}{mmol/L}. Puncak yang sesungguhnya jauh lebih rendah sebab penyerapannya
tersebar sepanjang sejam sementara pembuangannya menyingkirkan glukosanya ketika ia tiba, dan
hatinya mengambil sebagian besarnya pada lintasan pertamanya.
\textbf{3.} $\qty{75}{g}/\qty{60}{min} = \qty{1.25}{g/min}$; atas
\qty{15}{L}: yakni \qty{8.3}{mg/dL/min}, empat kali $u_{0}$.
\textbf{4.} Kelebihan mantapnya $u/a = 8.3/0.02 = \qty{417}{mg/dL}$, tetapan
waktunya $1/a = \qty{50}{min}$; sesudah sejam $417(1 -
\mathrm{e}^{-1.2}) = \qty{291}{mg/dL}$ di atas puasanya: yakni sekitar
\qty{380}{mg/dL}, \qty{21}{mmol/L}.
\textbf{5.} $L = 0.2\times 0.05/(0.02\times 0.1) = 5$; kelebihan mantapnya
$417/6 = \qty{69}{mg/dL}$, yakni enam kali lebih kecil daripada tanpa \omterm{def:b3:endocrinology:glucose}{insulin}.
\textbf{6.} Otaknya membakar \qty{120}{g} glukosa tiap hari, tidak menyimpannya
dan hanya sedikit dapat memakai yang lain: glukosa yang rendah memberi kebingungan dan koma dalam
hitungan menit. Glukosa yang tinggi mengglikasi protein lalu, bertahun-tahun, merusak
pembuluh kecil \omterm{def:b3:sensory-systems:retina}{retina}, ginjal, dan sarafnya serta nadi besarnya.
\textbf{7.} $\lambda^{2} + (a + \gamma)\lambda + (a\gamma + s\beta) = 0$:
jumlah diagonalnya $-0.12$, determinannya $0.002 + 0.01 = 0.012$, diskriminannya
$0.0144 - 0.048 = -0.0336$.
\textbf{8.} $\omega = \sqrt{0.0336}/2 = \qty{0.092}{min^{-1}}$; kalanya
\qty{68}{min}; waktu peluruhannya $2/0.12 = \qty{16.7}{min}$.
\textbf{9.} $\dot g(0) = 100(-0.06 + c\,\omega) = -2$, sehingga $c\,\omega =
0.04$, $c = 0.44$.
\textbf{10.} $g = 0$ ketika $\tan\omega t = -1/c = -2.29$, $\omega t =
\pi - 1.16 = 1.98$, $t = \qty{21.6}{min}$. Pada minimum pertamanya, dekat
\qty{32}{min}: $100\,\mathrm{e}^{-1.94}(\cos 2.97 + 0.44\sin 2.97) =
14.4\times(-0.91) \approx -\qty{13}{mg/dL}$, yakni glukosa \qty{77}{mg/dL}.
\textbf{11.} Pada $t = 0$, $\dot i = 100\beta > 0$: \omterm{def:b3:endocrinology:glucose}{insulinnya} naik sementara
$g > 0$, lalu memuncak ketika $\beta g = \gamma i$, yang menuntut $g$ masih
positif --- sehingga puncaknya datang sebelum glukosanya mencapai puasanya, dan sesudah
glukosanya mulai turun (ia turun sejak $t = 0$). Pada modelnya, puncaknya
dekat \qty{12}{min}.
\textbf{12.} Dengan masukan yang tersebar sepanjang sejam, puncaknya terjadi ketika
penyerapan dan pembuangannya seimbang, yang lebih rendah dan lebih lambat (30--60 menit);
kembalinya menunggu penyerapannya berakhir, karena itu dua jam; tembusan atas
\omterm{def:b3:endocrinology:glucose}{insulinnya} lebih kecil bagi sebuah masukan yang berangsur, sehingga lekuknya ringan.
\textbf{13.} $L = 2.5$. Yang sehat $g^{*} = (2/0.02)/6 = \qty{16.7}{mg/dL}$;
yang kebal $100/3.5 = \qty{28.6}{mg/dL}$: yakni sebuah kenaikan sekitar
\qty{12}{mg/dL} (\qty{0.7}{mmol/L}), dengan glukosa puasanya dekat
\qty{102}{mg/dL}, \qty{5.7}{mmol/L} --- yakni ambang glukosa puasa yang
terganggu.
\textbf{14.} $i^{*} = \beta g^{*}/\gamma$: yang sehat $8.3$, yang kebal
\qty{14.3}{mU/L}, nisbahnya $1.7$: yakni kekebalan \omterm{def:b3:endocrinology:glucose}{insulin} yang terimbangi,
dengan hiperinsulinemia dan glukosa yang nyaris normal.
\textbf{15.} $L = 0.1\times 0.01/0.002 = 0.5$; $g^{*} = 100/1.5 =
\qty{66.7}{mg/dL}$, yakni \qty{50}{mg/dL} di atas keadaan yang sehat:
\qty{2.8}{mmol/L}, dengan glukosa puasa \qty{7.8}{mmol/L} (\qty{140}{mg/dL}),
di atas ambang diabetesnya.
\textbf{16.} Diskriminannya $(a - \gamma)^{2} - 4s\beta = 0.0064 - 0.004 =
0.0024 > 0$: kembalinya monoton. $\lambda = -0.06 \pm 0.0245$: yakni $-0.0355$
dan $-0.0845$ tiap menit; tetapan waktu terlambatnya \qty{28}{min}, terhadap
\qty{16.7}{min} bagi peluruhan yang sehat: gelung yang gagal kembali lebih
lambat lalu tidak pernah menembus ke bawah.
\textbf{17.} Glukosa tinggi dengan \omterm{def:b3:endocrinology:glucose}{insulin} tinggi: yakni kekebalan dengan imbangan
sebagian --- jenis~2; nilai dua jamnya \qty{13}{mmol/L} melebihi
\qty{11.1}{}, sehingga diabetes, bukan sekadar toleransi yang terganggu.
\textbf{18.} Tanpa C-peptida berarti tanpa \omterm{def:b3:endocrinology:glucose}{insulin} endogen: yakni jenis~1. Tanpa
\omterm{def:b3:endocrinology:glucose}{insulin}, otot dan lemaknya tidak dapat mengambil glukosa, lemaknya dirombak menjadi
asam lemak dan keton, protein ototnya dikatabolismekan bagi
glukoneogenesisnya, dan glukosanya hilang di dalam kemihnya bersama kalorinya:
pasiennya kelaparan di tengah kelimpahan.
\textbf{19.} $5\times 9/5.5 = \qty{8.2}{\%}$.
\textbf{20.} $1.5\,\mathrm{e}^{0.46\times 3} = 1.5\times 3.97 =
\qty{6.0}{mU/L}$. T$_{4}$-nya di dalam jangkauannya, TSH-nya di atasnya: yakni
hipotiroidisme subklinis, dengan kelenjarnya gagal dan \omterm{def:b3:endocrinology:axes}{hipofisisnya} mengimbangi.
\textbf{21.} $1.5\,\mathrm{e}^{0.46\times 8} = 1.5\times 39.6 =
\qty{59}{mU/L}$. Sebuah T$_{4}$ rendah dengan sebuah TSH hanya $0.3$ berarti
\omterm{def:b3:endocrinology:axes}{hipofisisnya} tidak menanggapi: lesinya di pusat (\omterm{def:b3:endocrinology:axes}{hipofisis} atau
\omterm{def:b3:endocrinology:axes}{hipotalamus}), bukan di tiroidnya.
\textbf{22.} Satu paruh hidup: arasnya turun menjadi \qty{50}{\%}, dan
gejalanya merayap masuk sepanjang berminggu-minggu. Kortisol, dengan paruh hidup
\qty{80}{min}, lenyap dalam beberapa jam sesudah sebuah dosis yang terlewat, dan sebuah tekanan pada
hari itu tidak menjumpai pertahanan: sepekan tiroksin yang terlewat
tidak nyaman, sehari kortisol yang terlewat dapat membunuh.
\textbf{23.} T$_{4}$ tinggi (\omterm{def:b3:adaptive-immunity:antibody}{antibodinya} menggerakkan kelenjarnya tanpa memedulikan
umpan baliknya); TSH-nya tertekan oleh T$_{4}$ yang tinggi; hubungan log-linearnya
mendorongnya menjadi tak terdeteksi --- pada $T_{4} = 30$, $1.5\,\mathrm{e}^{-6.9}
= \qty{0.0015}{mU/L}$ --- sehingga sebuah TSH yang tak terdeteksi menjadi tanda pertamanya.
\textbf{24.} Aras sebuah \omterm{def:b3:endocrinology:hormone}{hormon} melaporkan neraca perintahnya dan
kelenjarnya: TSH tinggi dengan T$_{4}$ rendah adalah kelenjar yang gagal dan TSH tinggi dengan
T$_{4}$ tinggi adalah \omterm{def:b3:endocrinology:axes}{hipofisis} yang lepas kendali; \omterm{def:b3:endocrinology:glucose}{insulin} tinggi dengan glukosa tinggi adalah
tubuh yang kebal dan \omterm{def:b3:endocrinology:glucose}{insulin} rendah dengan glukosa tinggi adalah pulau yang musnah.
Bila dibaca sendirian, sebuah T$_{4}$ \qty{7}{pmol/L} atau sebuah \omterm{def:b3:endocrinology:glucose}{insulin}
\qty{30}{mU/L} tidak mengatakan di mana kesalahannya terletak; bila dibaca bersama
pemerintahnya, ia mengatakannya.
\textbf{25.} \omterm{thm:b3:endocrinology:loop}{Bati gelungnya} $L = 5$; kalanya \qty{68}{min} dan waktu peluruhannya
\qty{17}{min}; glukosa puasa gelung yang gagal \qty{7.8}{mmol/L}.
\end{solution}
